\begin{table}[htbp]\centering
\caption{Estimation-approach robustness (fig.\ D12)}
\begin{tabular}{lcccc}
\toprule
 & (1) & (2) & (3) & (4) \\
 & LP-DiD cells & LP-DiD within-event & Sale-level OLS & Sale-level OLS \\
\midrule
Pooled pre & +0.34 & -2.57 & -- & -- \\
 & (1.71) & (1.84) &  &  \\
Pooled post & +6.95*** & +5.10** & -- & -- \\
 & (2.03) & (2.15) &  &  \\
Post $-$ pre & +5.32*** & +3.16** & +7.09*** & +6.72*** \\
 & (1.40) & (1.43) & (1.58) & (1.32) \\
\midrule
Event$\times$ring FE & (diff.) & (diff.) & Y & Y \\
Event$\times$year FE & (diff.) & Y & Y & Y \\
Hedonic controls & N & N & N & Y \\
R$^2$ & 0.154 & 0.650 & 0.172 & 0.344 \\
Observations & 12,379 & 2,164 & 777,666 & 730,276 \\
\bottomrule
\end{tabular}
\begin{minipage}{0.95\textwidth}\vspace{2pt}\scriptsize The Post$-$pre row is the comparable differential across estimators. Column (1): cell-level LP-DiD, pooled never-treated controls; post$-$pre, R$^2$ and N from the single pre-mean-differenced regression (Dube et al.\ 2023). Column (2): the same PMD regression with event$\times$calendar-year effects absorbed (within-event identification: each near ring vs its own event's far ring); pooled pre/post are the corresponding lpdid coefficients. Columns (3)-(4): stacked sale-level OLS; the near$\times$post coefficient is itself the differential, SEs clustered by event. Long-differencing subsumes event$\times$ring FE in the LP-DiD columns. 100$\times$log points.\end{minipage}
\end{table}
